\documentclass[10pt,letterpaper]{article}
\usepackage{cornell-notes}

\title{Clause 26.07.28.05: Class Chunk View Inner Iterator}
\author{}
\date{}
\setDocTitle{Clause 26.07.28.05: Class Chunk View Inner Iterator}
\setDocSubtitle{C++ 2024}
\setDocOwner{C++ 2024 Cornell Notes}
\setDocDate{\today}
\setCornellCollection{C++ 2024 Cornell Notes}
\setCornellUnitType{Clause}
\setCornellUnitNumber{26.07.28.05}
\setCornellUnitTitle{Class Chunk View Inner Iterator}
\hypersetup{pdftitle={Clause 26.07.28.05: Class Chunk View Inner Iterator},pdfsubject={Cornell notes},pdfauthor={}}

\begin{document}
\maketitle
\begin{CornellOverviewBox}{Overview}
\textbf{Classification:} Standard library - Normative\\[2pt]
\textbf{Learning objective:} Understand the rules, terminology, and semantic consequences associated with Class chunk\_view::inner-iterator.
\end{CornellOverviewBox}\begin{CornellSummaryBox}{Summary}
{\sffamily\bfseries\color{Accent} Summary}\par\vspace{3pt}Section 26.7.28.5 focuses on Class chunk\_view::inner-iterator. Apply the specified interface together with its mandates, effects, result conditions, exception behavior, and complexity constraints. When applying it, preserve the distinction between mandatory requirements, recommendations, permissions, and behavior deliberately left undefined, unspecified, or implementation-defined.
\end{CornellSummaryBox}

\end{document}
